\chapter{Benchmark módszertan és a CSV-k felépítése}

\section{Benchmark futtatás lépései}

A WinForms alkalmazásban a \emph{Benchmark} fülön beállítható:

\begin{itemize}
  \item algoritmus: \textbf{CTR} vagy \textbf{GCM}
  \item kulcsméret: 128/192/256 bit
  \item adatméret: a mért bemenet mérete MB-ban (tipikusan 64 MB, néhol 256 MB)
  \item iterációk száma (a mért körök száma)
  \item warm-up: a mérés előtti bemelegítés (különösen OpenCL-nél fontos a cold-start miatt)
\end{itemize}

Minden iterációban azonos paraméterek mellett futtatok \textbf{Encrypt} és \textbf{Decrypt} méréseket,
mindhárom engine-re (amennyiben az adott kombináció támogatott).

\section{Mit jelentenek a mezők?}

A nyers CSV-k két részből állnak:

\begin{enumerate}
  \item \textbf{Meta sorok} \code{\#Kulcs,Érték} formában:
    környezet, gép, OS, build architektúra, OpenCL platform eszköz, stb.
  \item \textbf{Adatsorok} táblázatos formában: iteráció-szintű mérések.
\end{enumerate}

\subsection{Adatsor (táblázat) fontos oszlopai}

\begin{description}
  \item[Engine] \code{NativeCpu} / \code{OpenCl} / \code{ManagedAes}
  \item[Direction] \code{Encrypt} vagy \code{Decrypt}
  \item[ElapsedMilliseconds] a művelet nettó ideje \ms-ban
  \item[ThroughputMegabytesPerSecond] áteresztőképesség \mbps-ban
  \item[SpeedupVsSequentialNativeCpu] csak OpenCL-nél: gyorsulás a natív szekvenciális CPU-hoz képest
\end{description}

\textbf{Fontos:} a gyorsulást mindig a natív CPU baseline-hoz viszonyítom.
A .NET menedzselt AES-t csak \emph{referenciaként} tüntetem fel.

\section{Mérési környezetek (PC1--PC4)}

A CSV-kből a benchmark során automatikusan kiíratom a legfontosabb környezeti adatokat:
.NET verzió, Windows verzió/build, CPU, GPU, RAM, valamint OpenCL platform + device.

Az alábbi táblázat a CSV-kből generált, összefoglaló környezetlista (\code{generated/machines\_summary\_pgf.csv}).

\begin{table}[H]
  \centering
  \caption{A mérések fő környezeti jellemzői a CSV metaadatok alapján.}\label{tab:machines}
  \fittable{\pgfplotstabletypeset[
    col sep=comma,
    string type,
    columns={PC,BuildArchitecture,PowerSourceHu,WindowsBuild,ProcessorName,GpuName,OpenClPlatformName},
    columns/PC/.style={column name=PC, column type=c},
    columns/BuildArchitecture/.style={column name=Arch., column type=c},
    columns/PowerSourceHu/.style={column name=Táp, column type=c},
    columns/WindowsBuild/.style={column name=Windows build, column type=c},
    columns/ProcessorName/.style={column name=CPU, column type={>{\raggedright\arraybackslash}p{4.0cm}}},
    columns/GpuName/.style={column name=GPU, column type={>{\raggedright\arraybackslash}p{3.0cm}}},
    columns/OpenClPlatformName/.style={column name=OpenCL platform, column type={>{\raggedright\arraybackslash}p{2.6cm}}},
    every head row/.style={before row=\toprule, after row=\midrule},
    every last row/.style={after row=\bottomrule},
  ]{\GenPath/machines_summary_pgf.csv}}
\end{table}

\section{Megjegyzés a PC1 OpenCL x86 futásról}

A projekt során a \textbf{PC1} gépen az OpenCL-es \textbf{x86} build nem akart lefutni
(valószínűleg 32 bites OpenCL runtime/driver hiánya miatt).
Ez a dokumentumban külön megjegyzésként jelenik meg, és a PC1 esetében ezért csak x64-es OpenCL eredmények szerepelnek.

\cleardoublepage
